\subsection{Effective models over every o-minimal real closed field}

The all-field comparison makes the same finite PL codes describe
the entire definable concordance space in every chosen o-minimal
expansion.

\begin{corollary}\label{cor:effective-whole-definable-concordances}
Let $M$ be an o-minimal expansion of a real closed field $R$,
let $B$ be a closed definably compact definable manifold, and
choose a definable triangulation
$\tau:|K|_R\longrightarrow B$ with $K$ a finite abstract PL
manifold. The effective simplicial group
$G^{\mathrm{PL}}(K)$ of
Theorem~\ref{thm:effective-countable-concordance-models}
has a specified weak equivalence
\begin{equation}\label{eq:effective-whole-definable-model}
 G^{\mathrm{PL}}(K)\longrightarrow P_{\mathrm{def},M}(B).
\end{equation}
It is real-algebraic coefficient extension, followed by PL
inclusion and conjugation by $\tau\times\operatorname{id}_I$.
All families preserve the whole parameter, fix the bottom
pointwise, and have free top.

Uniformly from the finite simplex list of $K$ and $r\geq0$,
one can enumerate a complete countable presentation of
$\pi_rP_{\mathrm{def},M}(B)$. In degree zero this is the
full component group. In positive degrees the basepoint
is the identity, and translation gives the other components;
one may instead enumerate an integer relation matrix.
For the real realization $|K|_{\mathbb R}$, the same model
and presentations describe ordinary TOP concordances if
every component has dimension at most two or at least five.

The algorithm takes $K$ alone, with its manifold condition
as a promise. The triangulation specifies the comparison
map and need not be computable from a description of $B$.
No decision procedure for the theory of $M$ is an input.
The presentations are recursively enumerable; finiteness
and a terminating word problem are not asserted.
\end{corollary}

\begin{proof}
Let $F$ be the real algebraic field, embedded canonically
as an ordered subfield of $R$. The code construction gives
$G^{\mathrm{PL}}(K)=P_{\mathrm{PL},F}(|K|_F)$, and
Theorem~\ref{thm:effective-countable-concordance-models}
proves that coefficient extension to
$P_{\mathrm{PL},R}(|K|_R)$ is a weak equivalence.
Theorem~\ref{result-44} gives the next weak equivalence,
to $P_{\mathrm{def},M}(|K|_R)$.

Put $T=\tau\times\operatorname{id}_I$. On every parameter
simplex, conjugation sends a family $h_p$ to
$T h_pT^{-1}$. It preserves joint definability, continuity
of the total map and its inverse, the fixed bottom and
the free top. It commutes with every face and degeneracy
because $T$ is independent of the parameter. Conjugation
by $T^{-1}$ is its inverse. It is therefore an isomorphism
of definable simplicial groups, proving
\eqref{eq:effective-whole-definable-model} by composition.

For completeness, the presentation constructed earlier
uses the decidable cycle set $Z_r$ and Moore subgroup
$N_{r+1}$ in the coded group $G$. Its relations are
\[
 x_1=1,\qquad x_ax_b=x_{ab}\quad(a,b\in Z_r),
 \qquad x_{d_0u}=1\quad(u\in N_{r+1}).
\]
The multiplication relations present all of $Z_r$ and
the boundary relations quotient by precisely $d_0N_{r+1}$.
The first-face Moore formula established in the preceding
theorem identifies this quotient with $\pi_r|G|$, including
$r=0$. Enumerating finite strings, and pairs of strings,
up to each successive length enumerates every required
generator and relation. Every stage uses only terminating
tests on the finite codes for $K$. The comparison with
$M$ changes the interpretation of this presentation,
not the algorithm. Positive-degree groups are abelian,
so their multiplication and boundary relations give the
stated integer matrix as before.

For a real base, Theorem~\ref{result-4} gives the ordinary
PL--TOP comparison in dimensions at least five, and
Theorem~\ref{thm:low-dimensional-real-concordance-comparison}
gives it in dimensions at most two. A bottom-fixed
concordance preserves every component of $|K|\times I$:
each component meets its own fixed bottom. The entire
group is consequently the finite product of the groups
of those components, and the comparisons apply to that
product, including when the component dimensions differ.
For $K=\varnothing$, all groups are the one-element group.
\end{proof}
